% Chapter 11 back matter: physical PDF p. 521 (printed p. 498).
% Inventory: four problems and nine references; no figures or footnotes.

\chapterbackmatter{Problems}

\begin{enumerate}
\item Calculate the contributions to the vacuum polarization function
$\pi(q^2)$ and to $Z_3$ of one-loop graphs involving a charged spinless
particle of mass $m_s$. What effect does this have on the energy shift of the
$2s$ state of hydrogen, if $m_s \gg Z\alpha m_e$?

\item Suppose that a neutral scalar field $\phi$ of mass $m_\phi$ has an
interaction $g\phi\bar\psi\psi$ with the electron field. To one-loop order,
what effect does this have on the magnetic moment of the electron? On $Z_2$?

\item Consider a neutral scalar field $\phi$ with mass $m_\phi$ and
self-interaction $g\phi^3/6$. To one-loop order, calculate the $S$-matrix
element for scalar-scalar scattering.

\item To one-loop order, calculate the effect of the neutral scalar field of
Problem 2 on the mass shift $\delta m_e$ of the electron.
\end{enumerate}

\chapterbackmatter{References}

\begin{enumerate}
\item\label{ref:11.1} R. P. Feynman, \emph{Phys. Rev.} \textbf{76}, 769 (1949).

\item\label{ref:11.2} G. C. Wick, \emph{Phys. Rev.} \textbf{80}, 268 (1950).

\item\label{ref:11.3} G. 't Hooft and M. Veltman, \emph{Nucl. Phys.} \textbf{B44}, 189
(1972).

\item\label{ref:11.4} E. A. Uehling, \emph{Phys. Rev.} \textbf{48}, 55 (1935). The one-loop
function $\pi(q^2)$ was first given for $q^2 \ne 0$ by J. Schwinger,
\emph{Phys. Rev.} \textbf{75}, 651 (1949).

\item\label{ref:11.5} J. Schwinger, \emph{Phys. Rev.} \textbf{73}, 416 (1948).

\item\label{ref:11.6} This is calculated (including terms that vanish for $m_e \ll m_\mu$)
by H. Suura and E. Wichmann, \emph{Phys. Rev.} \textbf{105}, 1930 (1967);
A. Petermann, \emph{Phys. Rev.} \textbf{105}, 1931 (1957); H. H. Elend,
\emph{Phys. Lett.} \textbf{20}, 682 (1966); \textbf{21}, 720 (1966);
G. W. Erickson and H. H. T. Liu, UCD-CNL-81 report (1968).

\item\label{ref:11.7} J. Bailey \emph{et al.} (CERN--Mainz--Daresbury Collaboration),
\emph{Nucl. Phys.} \textbf{B150}, 1 (1979). These experiments are done by
observing the precession of the muon spin in a storage ring.

\item\label{ref:11.8} W. Pauli and F. Villars, \emph{Rev. Mod. Phys.} \textbf{21}, 434
(1949). Also see J. Rayski, \emph{Phys. Rev.} \textbf{75}, 1961 (1949).

\item\label{ref:11.9} See, e.g., A. I. Miller, \emph{Theory of Relativity --- Emergence
(1905) and Early Interpretation (1905--1911)} (Addison-Wesley, Reading, MA,
1981): Chapter 1.
\end{enumerate}
